\subsection{The Distributional Hypothesis}
\label{sec:distributional-hypothesis}
\slidesrc{Distributional semantics: 3--16; Semantics (long version): 3--5}

\subsubsection{Meaning and distribution}
\slidesrc{Distributional semantics: 4--7}

Four famous quotations, each followed by what it means for us:

\begin{itemize}
  \item ``\emph{Die Bedeutung eines Wortes liegt in seinem Gebrauch}'' (``the meaning of
        a word is its use'', Ludwig Wittgenstein).\\
        $\Rightarrow$ \textbf{meaning = use = distribution in language}.
  \item ``\emph{You shall know a word by the company it keeps!}'' (J.\,R.~Firth, 1957).\\
        $\Rightarrow$ \textbf{distribution = collocations} = habitual word combinations
        (Section~\ref{sec:corpora}).
  \item \textbf{Distributional hypothesis}: difference of meaning correlates with
        difference of distribution (Zellig Harris, 1954).\\
        $\Rightarrow$ we can measure a \textbf{semantic distance} between words.
  \item ``\emph{What people know when they say that they know a word is not how to recite
        its dictionary definition -- they know how to use it [\ldots] in everyday
        discourse}'' (George Miller, 1986).
\end{itemize}

\begin{nbdef}{Distributional hypothesis}
Words that occur in similar contexts tend to have similar meanings. Conversely, the more
the contexts of two words differ, the more their meanings differ.
\end{nbdef}

\subsubsection{What is the meaning of ``bardiwac''?}
\slidesrc{Distributional semantics: 8--16}

You have never seen the word \emph{bardiwac}. Read these sentences, one at a time:
\begin{itemize}[nosep]
  \item He handed her her glass of bardiwac.
  \item Beef dishes are made to complement the bardiwacs.
  \item Nigel staggered to his feet, face flushed from too much bardiwac.
  \item Malbec, one of the lesser-known bardiwac grapes, responds well to Australia's
        sunshine.
  \item I dined off bread and cheese and this excellent bardiwac.
  \item The drinks were delicious: blood-red bardiwac as well as light, sweet Rhenish.
\end{itemize}

\begin{nbintuition}[what you just did]
After a few sentences you know that \textbf{bardiwac is a heavy red alcoholic beverage
made from grapes}, i.e.\ a red wine (the word is invented; it stands for
\emph{claret}). Nobody gave you a definition: you inferred it from the
\emph{company the word keeps}: \emph{glass of}, \emph{too much}, \emph{grapes},
\emph{drinks}, \emph{blood-red}, \emph{beef dishes}. That is the distributional hypothesis
in action, and it is how we learn most of our vocabulary.
\end{nbintuition}

The examples are hand-picked and edited, of course. But in a corpus like the British
National Corpus (BNC) there is at least as much relevant information. The slides show a
\emph{word sketch} of the real word (BNC frequency 230): grouped by grammatical relation,
its most typical companions are verbs like \emph{uncork}, \emph{gulp}, \emph{sip},
\emph{drink}, \emph{swig}; coordinated nouns like \emph{burgundy}, \emph{Sancerre},
\emph{scampi}, \emph{ruby}; containers like \emph{bottle}, \emph{glass}, \emph{jug},
\emph{goblet}; modifiers like \emph{full-bodied}, \emph{blood-red}, \emph{ready-to-drink}.
A computer can collect these statistics automatically.
